\begin{textsample}{POS Dim 2 – human – Score 7.00 – es/es\_ef\_anos\_iniciais\_ensino\_religioso\_1\_p9.txt}  \label{ex:f2_pos_009}
No Ensino Fundamental, o Ensino Religioso adota a pesquisa e o diálogo como princípios mediadores e articuladores dos processos de observação, identificação, análise, apropriação e ressignificação de \textbf{saberes}, visando o \textbf{desenvolvimento} de competências específicas.
Dessa maneira, \textbf{busca} problematizar representações sociais preconceituosas sobre o outro, com o intuito de combater a intolerância, a discriminação e a exclusão.
Por isso, a interculturalidade e a ética da alteridade constituem fundamentos teóricos e pedagógicos do Ensino Religioso, porque favorecem o reconhecimento e respeito às histórias, memórias, crenças, convicções e valores de diferentes culturas, tradições religiosas e filosofias de vida.
O Ensino Religioso \textbf{busca} construir, por \textbf{meio} do estudo dos conhecimentos religiosos e das filosofias de vida, atitudes de reconhecimento e respeito às alteridades.
Trata -se de um espaço de aprendizagens, experiências pedagógicas, intercâmbios e diálogos permanentes, que visam o acolhimento das identidades culturais, religiosas ou não, na perspectiva da interculturalidade, direitos \textbf{humanos} e cultura da paz.
Tais finalidades se articulam aos elementos da formação integral dos estudantes, na medida em que fomentam a aprendizagem da convivência democrática e cidadã, princípio básico à vida em \textbf{sociedade}.

% matched lemmas: buscar, desenvolvimento, humano, meio, saber, sociedade
\end{textsample}
